A binary search tree answers a query in time proportional to its height, and
nothing in the search tree property forces that height to stay small. Keys that
arrive in sorted order build a tree of $n$ nodes and height $n-1$, which is a
linked list wearing the interface of a tree. The AVL tree, published by
Adelson-Velsky and Landis in 1962, removes that failure by requiring the two
subtrees of every node to differ in height by at most one, and by restoring the
requirement after every update.

We give a self-contained account of the structure. We derive the height bound
$h < 1.4404\log_2(n+2) - 1.3277$ from the recurrence that counts the nodes of a
minimum-size tree, we prove that one insertion needs at most one rebalancing
step while one deletion may need a rebalancing step at every level, and we give
the algorithms in full. We then measure an implementation on inputs of up to one
million keys. The measurements confirm the analysis and sharpen it: the height
bound holds but is not tight, and although the worst case allows a deletion to
rotate at every level, deletions average $0.27$ rebalancing events apiece and do
not grow with $n$.
